\section{Matching the physical parent and integrating the conditional mixture}
\label{sec:parent-conditioned-repair}

The aggregate-moment and component-mixture investigations raise a more specific
question than whether a flexible guide can fit a useful distribution. Which
distribution does the flow actually transport? A Gaussian fitted to the whole
predictive cloud describes uncertainty across parents. A child generated from
one physical parent starts with the process covariance around that parent's
transition mean. These distributions generally have different means and
covariances. Applying a flow designed for the first to samples from the second
can produce a poor proposal even when its density and importance correction are
implemented correctly. The correction preserves the stated importance law;
it does not make a poor proposal efficient.

This section derives a repair in the SIR setting, where the nonlinear dynamics
are followed by a linear Gaussian observation. It connects the conditional
proposal already introduced in Section~\ref{sec:sg-physical-repair} to exact
integration of the current finite predictive mixture. The proposal itself is
an existing option, not a newly discovered algorithm. The new investigation
separates its remaining ancestor error from the extra error introduced by
representing each conditional Gaussian with one child and then resetting the
cloud. It also specifies what may be changed in the KDM construction without
silently changing either the importance law or the differentiated scalar.

\subsection{Why a covariance mismatch can become a mean error}

Consider a scalar observation $y=x+\varepsilon$, with
$\varepsilon\sim N(0,R)$. Suppose the flow is designed for the prior $N(m,P)$,
but its input really follows $N(a,Q)$. The continuous affine Gaussian flow that
transports its intended prior to its posterior has terminal map
\begin{equation}
 F_{m,P}(x)=m+\frac{P}{P+R}(y-m)
              +\sqrt{\frac{R}{P+R}}(x-m).
 \label{eq:parent-flow-map}
\end{equation}
Indeed, when $x\sim N(m,P)$, this map has the posterior mean
$m+P(y-m)/(P+R)$ and variance $PR/(P+R)$. On the actual physical input,
however, the mean and variance are
\begin{align}
 E[F_{m,P}(X)]&=m+\frac{P}{P+R}(y-m)
                 +\sqrt{\frac{R}{P+R}}(a-m),\\
 \operatorname{Var}[F_{m,P}(X)]&=\frac{RQ}{P+R}.
 \label{eq:parent-mismatch}
\end{align}
The desired parent-conditional Gaussian instead has mean
$a+Q(y-a)/(Q+R)$ and variance $RQ/(Q+R)$. Setting $m=a$ removes the
centering mismatch but leaves the wrong gain and variance when $P\ne Q$.
Setting $P=Q$ is also insufficient if $m\ne a$: with
$c=\sqrt{R/(Q+R)}$, the remaining mean error is
$(c^2-c)(m-a)$. Matching both quantities makes the continuous linear-Gaussian
map exact. A finitely discretized flow still requires its own integration
error check.

Consequently, a mean-dominated discrepancy does not show that covariance
repair is irrelevant. The covariance controls both the gain and contraction,
and therefore the output mean. For example, take $a=m=0$, $Q=1$, $R=100$,
and $y=10$. The physical conditional mean is $10/101\simeq0.099$, with
variance $100/101\simeq0.990$. A flow designed with $P=100$ instead produces
mean $5$ and variance $0.5$ on those same physical draws. Nothing in this
example requires a failed Cholesky factorization or extreme floating-point
rounding: it is a mismatch between the assumed and actual proposal input.

The R2 frozen-input contraction diagnostic supports this mechanism in a
limited, descriptive sense. Replacing the guide covariance by
$P_\gamma=Q+\gamma(P_{\rm UKF}-Q)$, while preserving its mean and recomputing
the proposal densities, reduced the observed mean conditional Gaussian KL
from about $81.14$ at $\gamma=1$ to $0.002$ at $\gamma=0$ on the inspected
stages. Those stages do not establish a general ranking or show that the
remaining mean mismatch is always negligible. They motivate a cleaner test:
condition the actual parent distribution directly.

\subsection{The exact conditional law of the current finite mixture}

Fix a time step and parameter $\theta$. The incoming approximation consists
of states $x_i$ and normalized positive weights $w_i$, $i=1,\ldots,N$.
Let $a_i=f_\theta(x_i)$ be the nonlinear transition mean. In this SIR model,
the transition is Gaussian conditional on the parent, and the observation is
linear:
\begin{equation}
 X\mid i\sim N(a_i,Q),\qquad Y\mid X\sim N(HX,R),\qquad
 p_N(x)=\sum_iw_i\,\varphi(x;a_i,Q),
 \label{eq:parent-finite-prior}
\end{equation}
where $\varphi(\cdot;m,V)$ denotes a normalized Gaussian density. The state has
18 coordinates; $H$ selects the nine infected coordinates. The physical
covariances in the present implementation are $Q=I_{18}$ and
$R=100\exp(2\theta_2)I_9$, where $\theta_2$ is the observation-noise
coordinate in the code's zero-based indexing. Thus $Q$ is not the covariance
of the whole predictive cloud and should not be replaced by that covariance
when defining the conditional physical transition.

Define the innovation covariance and Kalman gain
\begin{equation}
 S=HQH^\top+R,\qquad K=QH^\top S^{-1}.
\end{equation}
Completing the square in the product of transition and observation densities
gives
\begin{align}
 m_i&=a_i+K(y-Ha_i),& C&=Q-KHQ,\\
 e_i&=\varphi(y;Ha_i,S),&
 \varphi(x;a_i,Q)\varphi(y;Hx,R)
   &=e_i\varphi(x;m_i,C).
 \label{eq:parent-product-identity}
\end{align}
The last equality includes its normalizing constant: integrating both sides
over $x$ gives the predictive likelihood $e_i$ for parent $i$. Summing with
the incoming weights then yields
\begin{equation}
 Z_N=\sum_iw_i e_i,\qquad
 \alpha_i=\frac{w_i e_i}{Z_N},\qquad
 \pi_N(x\mid y)=\sum_i\alpha_i\varphi(x;m_i,C).
 \label{eq:parent-exact-posterior}
\end{equation}
This is exact conditioning of the finite mixture in
\eqref{eq:parent-finite-prior}. The incoming states remain an approximation
to the previous filtering distribution. Therefore $Z_N$ is not generally the
true model predictive likelihood, and $\pi_N$ is not an exact solution of the
nonlinear recursive filter. Direct conditioning removes one current-step
approximation while leaving the previous approximation visible.

\subsection{What changes in sampling, and what does not}

Use $q_i(x)=\varphi(x;m_i,C)$ for the child associated with parent $i$.
Equation~\eqref{eq:parent-product-identity} gives the pointwise cancellation
\begin{equation}
 \frac{\varphi(x;a_i,Q)\varphi(y;Hx,R)}{q_i(x)}=e_i.
 \label{eq:parent-cancellation}
\end{equation}
With one child drawn from each parent, the unnormalized child weight is
$w_i e_i$. There is no child-draw variability in this ratio. The parent
evidences can nevertheless differ enormously, so the normalized weights
$\alpha_i$ can still concentrate. This repair cannot recover a missing region
of the incoming parent cloud.

The allocation law matters if one changes that one-child-per-parent design.
If a parent is sampled with probability $b_i$ and its child is sampled from
$q_i$, the joint importance ratio is
\begin{equation}
 \frac{w_i\varphi(x;a_i,Q)\varphi(y;Hx,R)}{b_iq_i(x)}
       =\frac{w_ie_i}{b_i}.
 \label{eq:parent-allocation-ratio}
\end{equation}
Choosing $b_i=\alpha_i$ makes that joint ratio constant, but a categorical
choice based on parameter-dependent $\alpha_i$ makes the realized finite
program piecewise discontinuous in $\theta$. Between switching boundaries,
the selected branch has an ordinary pathwise derivative; at a jump, that
derivative does not describe passage across the boundary. It therefore does
not establish the smooth finite target required by the present HMC design.
Differentiating an expectation over the categorical choice is a separate
question and must also account for the changing allocation probabilities. This
allocation is now an implemented proposal default, as detailed in
Section~\ref{sec:observation-weighted-default}; its HMC smoothness limitation
remains. Smooth deterministic transport is a different route;
it needs its own finite-map derivative and reset analysis.

Likewise, replacing an ancestor density by a marginal mixture density requires
a corresponding derivation of the sampling measure. A global KDM density
cannot simply be inserted into the denominator of a parent-specific ratio
while keeping an incompatible numerator or allocation weight. In the first
repair experiment, parent association and the existing ancestor importance
law remain fixed. The shared physical-conditional proposal already computes
the original transition, observation and proposal densities at the same
child; the algebraic cancellation is checked rather than assumed.

There is no double assimilation here. The observation enters the proposal,
and the original target numerator still contains the observation likelihood;
the proposal density in the denominator makes
\eqref{eq:parent-cancellation} hold. Applying another observation update to the
conditional Gaussian without adjusting the proposal law would define a
different calculation.

\subsection{Integrate before blaming the child representation}

The finite conditional mixture permits exact first and second moments without
drawing a child. Its mean, raw second moment and covariance are
\begin{align}
 \mu_N&=\sum_i\alpha_i m_i,\\
 M_{2,N}&=\sum_i\alpha_i(C+m_im_i^\top),\\
 V_N&=\sum_i\alpha_i\{C+(m_i-\mu_N)(m_i-\mu_N)^\top\}.
 \label{eq:parent-integrated-moments}
\end{align}
The covariance expression separates within-parent uncertainty $C$ from
between-parent uncertainty. Keeping only the weighted conditional means would
discard $C$ and underestimate the uncertainty even when all parents coincide.
If they do coincide, \eqref{eq:parent-integrated-moments} reduces to the
ordinary Gaussian conditional moments, a useful implementation check.

For the conditional proposal, the weights $\alpha_i$ are independent of the
child draws. Conditional on the incoming cloud, the expectation of
$\sum_i\alpha_i\phi(X_i)$ equals the corresponding integrated expectation
for any integrable $\phi$. Thus integration eliminates that conditional
sampling variance for first and quadratic features. With independent children,
the conditional variance of the sampled feature estimator is
$\sum_i\alpha_i^2\operatorname{Var}_{q_i}[\phi(X)]$.

There is also a precise squared-error statement. Write
$\widehat g=\sum_i\alpha_i\phi(X_i)$ and
$g_N=\sum_i\alpha_i E_{q_i}[\phi(X)]$. Conditional on the incoming cloud
and observation, let $g_\star$ be a fixed reference feature vector and let
$W$ be a fixed positive semidefinite loss matrix. Expanding the square and
using $E[\widehat g-g_N]=0$ gives
\begin{equation}
 E\!\left[(\widehat g-g_\star)^\top W(\widehat g-g_\star)\right]
 =(g_N-g_\star)^\top W(g_N-g_\star)
 +\operatorname{tr}\!\left\{W\sum_i\alpha_i^2
       \operatorname{Cov}_{q_i}[\phi(X)]\right\}.
 \label{eq:parent-integrated-risk}
\end{equation}
Thus integration reduces expected squared error relative to these sampled,
pre-reset conditional-proposal features, even if $g_\star$ is the true
filtering expectation. The same identity applies to an independently generated
latent-state label after conditioning on that label. It leaves the first term,
the incoming approximation's error, untouched. This expectation over fresh
child draws does not order individual realized losses, arbitrary
self-normalized importance estimators, corrected clouds, or two recursive
algorithms with different incoming clouds. The identity also concerns exact
arithmetic; finite-precision departures from the cancellation must be measured.

The comparison becomes informative when all readouts are retained: integrated,
weighted children, reset cloud and corrected cloud. If integrated moments are
already poor on a fixed incoming cloud, changing the number of child draws
cannot fix that conditional target; investigate incoming coverage and earlier
propagation. If integration is accurate but sampled moments are poor, inspect
allocation and within-component representation. If the discrepancy first
appears after reset or correction, inspect those transformations. These are
explanatory comparisons with the same incoming approximation, not competing
definitions of the model posterior.

The canonical Contract-E reset restores the moments of its weighted source
cloud. Feeding $\mu_N$ and $V_N$ into that reset instead would change its
target and therefore require a distinct reset contract and a new total
derivative. The initial implementation consequently records integrated moments
without using them to alter the next cloud. An integrated final readout can
improve the reported estimator while leaving next-step ancestor coverage
unchanged; both effects must be measured separately.

\subsection{Differentiate the same conditioned mixture}

For a parameter direction, write a dot for the total directional derivative.
The parent means depend on both the transition parameter and the incoming
states: $\dot a_i=\partial_\theta f_\theta(x_i)\dot\theta+
D_xf_\theta(x_i)\dot x_i$. Incoming weights, previous resets and corrections
also depend on the parameter. Holding these quantities fixed would answer a
different derivative question.

To allow parameter-dependent $H,Q,R$, differentiate the innovation covariance
and gain explicitly:
\begin{align}
 \dot S&=\dot H QH^\top+H\dot QH^\top+HQ\dot H^\top+\dot R,\\
 \dot K&=(\dot QH^\top+Q\dot H^\top-K\dot S)S^{-1},\\
 r_i&=y-Ha_i,& \dot r_i&=-\dot H a_i-H\dot a_i,\\
 \dot m_i&=\dot a_i+\dot K r_i+K\dot r_i,\\
 \dot C&=\dot Q-\dot KHQ-K\dot H Q-KH\dot Q.
 \label{eq:parent-condition-tangents}
\end{align}
The observation $y$ is held fixed during inference. Matrix solves implement
these expressions without forming an explicit inverse. With
$v_i=S^{-1}r_i$, the Gaussian log-evidence derivative is
\begin{equation}
 \dot{\log e_i}=-v_i^\top\dot r_i+
 \tfrac12v_i^\top\dot S v_i-\tfrac12\operatorname{tr}(S^{-1}\dot S).
\end{equation}
For normalized incoming weights, put
$u_i=\dot{\log w_i}+\dot{\log e_i}$. Then
\begin{equation}
 \dot{\log Z_N}=\sum_i\alpha_i u_i,\qquad
 \dot\alpha_i=\alpha_i\left(u_i-\sum_j\alpha_j u_j\right).
 \label{eq:parent-weight-tangents}
\end{equation}
When the input consists of unnormalized log weights $\ell_i$, replace
$\dot{\log w_i}$ by
$\dot\ell_i-\sum_jw_j\dot\ell_j$. This normalization term is necessary even
when all displayed weights sum numerically to one.

Finally, let $B_i=C+m_im_i^\top$. The readout derivatives are
\begin{align}
 \dot\mu_N&=\sum_i(\dot\alpha_i m_i+\alpha_i\dot m_i),\\
 \dot M_{2,N}&=\sum_i\left[\dot\alpha_i B_i+
 \alpha_i(\dot C+\dot m_i m_i^\top+m_i\dot m_i^\top)\right].
 \label{eq:parent-readout-tangents}
\end{align}
These derivatives include mixture weights as well as component moments.
They can be evaluated with the analytical recursive tangents already carried
by the shared executor. Finite differences test the total recursive program,
including moving initial states and all previous resets. Agreement certifies
that finite-program derivative on the tested cases, not agreement with the
underlying model score or with a Zhao--Cui adapted-filter reference.

\subsection{Consequences for a useful KDM repair}

For a Gaussian mixture whose components approximate the predictive
distribution, condition each component before combining it. Component $k$
with predictive mean $a_k$ and covariance $Q_k$ has its own $S_k,K_k,m_k,C_k$
and evidence $e_k$; update its weight proportionally to $w_ke_k$. Aggregating
first loses this component-specific observation response. But merely creating
more components cannot guarantee an improvement. For instance, the symmetric
scalar mixture with means $\pm b$ and common variance $P-b^2$ matches the
mean and variance of $N(0,P)$, while its fourth moment is $3P^2-2b^4$.
Moment matching has not preserved the Gaussian law.

The physical-parent repair and a useful learned or UKF-updated KDM therefore
answer different parts of the problem. Matching $a_i,Q$ repairs the flow's
local input mismatch. Integrating its conditional mixture removes current-step
child noise from the chosen readout. Improving later incoming coverage may
require a better component representation or smooth redistribution, but that
step must preserve a derived sampling law and include the corresponding
analytical derivatives. The first execution tests the existing conditional
proposal with unchanged reset semantics and reports all four moment readouts.
Only after that comparison identifies the remaining error is a more elaborate
component or allocation mechanism justified.

The initial independent Gaussian-product check passed in FP32 and FP64; its
FP64 full moment/evidence tangent discrepancy was below $6\times10^{-10}$
on a three-parent fixture. That fixture held $H$ and $Q$ fixed and did not test
the recursive endpoint. Its purpose is to check
\eqref{eq:parent-product-identity} and the moving-weight/readout algebra.
Recursive, GPU and paired estimation results belong to the subsequent
execution record; the fixture alone establishes no efficacy ranking.

\subsection{What the first recursive comparison establishes}

The first recursive experiment uses the 18-state SIR model, five observations,
1,008 particles, 16 newly simulated datasets and two particle replicates per
dataset, at the nominal three-dimensional parameter vector. It compares the
original LEDH baseline, the existing aggregate-moment candidate with
$\gamma=0$, and the physical-parent conditional proposal. The latter uses one
Gaussian child per parent and retains the complete covariance, Contract-E
reset and correction lifecycle. All three routes run through the same
analytical executor. Integration supplies an additional readout; it does not
change the cloud propagated to the next observation.

The loss is the equal-weight average of squared errors in two feature blocks:
18 standardized first moments and the 171 distinct standardized quadratic
moments, evaluated against the simulated terminal state. Feature scales were
frozen using a separate 64-path calibration. Consequently, the numbers below
measure a latent-label prediction loss, not distance to the exact filtering
moments or likelihood. UKF, bootstrap filtering and the UKF prediction before
the terminal observation provide simple comparators.

The numerical-reference run uses FP32 particles and densities, disables TF32,
and retains the previously selected FP64 output standardization within the
correction. It is therefore a mixed-precision run, not evidence for an
all-FP32 implementation. Every one of its 96 full-filter cells passes the
recorded validity checks. A separate CPU/FP64 two-observation fixture checks
the total recursive derivatives of the likelihood and integrated and corrected
moments by central differences at two step sizes, including dependence through
the initial cloud and earlier resets.

\begin{center}
\begin{tabular}{lr}
\hline
Method and final readout & Mean standardized feature loss\\
\hline
UKF & 0.566871\\
Bootstrap particle filter & 0.602845\\
Aggregate-moment candidate, $\gamma=0$, corrected & 0.605266\\
Parent-conditioned proposal, weighted children & 0.606323\\
Parent-conditioned proposal, reset cloud & 0.606323\\
Parent-conditioned proposal, corrected cloud & 0.606323\\
Parent-conditioned proposal, integrated & 0.607334\\
UKF prediction before the terminal observation & 0.656055\\
Original LEDH baseline, corrected & 0.719388\\
\hline
\end{tabular}
\end{center}

These are descriptive means over datasets after averaging the two particle
replicates. The parent-conditioned corrected loss minus the $\gamma=0$
corrected loss is $0.001057$, with a dataset-paired standard error of
$0.000949$. This screen does not establish an accuracy improvement over that
existing candidate. The integrated loss minus the weighted-child loss is
$0.001011$, with a paired standard error of $0.000396$. That realized
difference does not contradict \eqref{eq:parent-integrated-risk}: the identity
averages over fresh child draws conditional on each incoming cloud, whereas
the table uses two particular draws and one latent label per dataset.
The screen was specified as descriptive and selects no statistically
supported winner.

The comparison with simple methods is consequential. The parent-conditioned
integrated loss exceeds UKF by $0.040464$ overall, with a paired standard
error of $0.016877$, and its observed loss is higher in each of eight
predeclared low/high regime bins. Those bins split innovation, predictive
spread, signal-to-noise ratio and distance from the boundary at thresholds
frozen on earlier calibration data; they overlap and contain between five and
eleven datasets. This is a promotion veto under the experiment's conditional
heuristic screen, not a claim of a population ranking or eight independent
replications.

The four readouts localize a different question. Across the 32
parent-conditioned runs and all five times, the largest coordinatewise
absolute difference between sampled and integrated means is $0.4181$; for
covariance entries it is $5.0501$. The corresponding reset-minus-sample
maxima are $9.89\times10^{-4}$ and $3.16\times10^{-4}$, while
correction-minus-reset maxima are $7.60\times10^{-5}$ and
$9.09\times10^{-5}$. Thus reset and correction approximately preserve the
moments of their sampled source in this reference run. These discrepancies
measure agreement with that source, not error relative to the true filter.
In particular, integrating the same incoming approximation has not closed
its observed gap to UKF. Incoming coverage, the finite-parent representation
and the chosen proposal family remain explanations to distinguish.

TF32 requires a separate numerical conclusion. Repeating the screen with
identical saved data and random inputs leaves all 32 $\gamma=0$ cells and
all 32 parent-conditioned cells valid, but rejects all 32 original-baseline
cells: 31 first fail at reset and one first fails at correction. The failed
baseline cells have effective sample sizes between 1 and 3.03 at the failing
time. Reset row and column residuals remain below their declared thresholds,
so the recorded failures are not marginal-balance failures. The records locate
the failing stage; they do not identify an individual failing factorization.
An accuracy comparison against that baseline in TF32 would therefore be invalid.

The parent-conditioned cancellation is also more accurate without TF32.
Across the reference screen, the largest absolute discrepancies between
sampled importance and integrated formulas are $6.08\times10^{-7}$ for a
normalized weight, $3.06\times10^{-5}$ for the cumulative log evidence and
$6.49\times10^{-5}$ for its tested directional derivative. With TF32 these
grow to $6.61\times10^{-4}$, $1.82\times10^{-3}$ and $1.43\times10^{-2}$,
respectively. Internal validity flags consequently do not imply precision
equivalence or HMC readiness.

The physical-parent mismatch now has a derived repair and a tested recursive
implementation. Its integrated moments are exact for the stated finite
conditional mixture, up to numerical error. Neither statement establishes
that this incoming mixture adequately represents the nonlinear filtering
distribution. The next scientific question is whether improving that
representation or smooth allocation closes the remaining error; it is
different from repeating the conditional Gaussian algebra. No Zhao--Cui
adapted-filter oracle or HMC run was part of this experiment.

The complete commands, frozen controls, paired regime tables, failure census
and source snapshots are recorded in the repository note
\texttt{ledh-parent-conditioned-result-20261026.md}. The corrected no-TF32
and TF32 runs are identified there by timestamps
\texttt{20261026T075252756468Z} and \texttt{20261026T075603594553Z}.
